% Storage formats and annotation codes: the byte layouts of the signal formats, the FLAC
% formats, the MIT and AHA annotation file formats, and the tables of annotation codes, as
% implemented by python/src/wfdbx/_signal.py, _flac.py, _annot.py and _ecgcodes.py (and
% identically by the other implementations).
\chapter{Storage formats and annotation codes}
\label{chap:formats}

This chapter is the reference of the binary files of a record: the storage formats of the
signal files, named by the format numbers of the header (\cref{chap:records}), the layout
of the annotation files, and the annotation codes with their mnemonics, descriptions and
mapping tables. These are the formats of the WFDB library \wfdbversion{}; the files that
\wfdbx{} writes are read by the C library, and those of the C library by \wfdbx{}, in all
seven implementations.

\section{Signal storage formats}
\label{sec:formats-signal}

A signal file holds the samples of one \emph{signal group}: the signals of the header that
name the same file. The samples are stored frame by frame: for each frame, the samples of
the first signal of the group (one, or \textit{n} for a signal with \textit{n} samples per
frame, written \code{16x}\textit{n} in the header), then those of the second signal, and so
on. The formats that pack several samples into a group of bytes (212, 310 and 311) pack
consecutive samples of this sequence, whatever signals they belong to. The header can also
give a byte offset (\code{16+44}: the samples begin after 44 bytes of prolog, which the
library skips) and a skew for each signal (\cref{chap:records}).

Each format represents a range of integer values (ADC units). The most negative value of a
format is reserved: it marks an \emph{invalid} sample (a sample that is missing or was not
recorded). \code{putvec} writes the invalid sample \code{WFDB\_INVALID\_SAMPLE} ($-32768$) as
that value, and \code{getvec} reads that value back as \code{WFDB\_INVALID\_SAMPLE} (or, in
the \code{WFDB\_GVPAD} mode of \code{setgvmode}, as the previous valid sample of the signal).
A sample outside the range of the format is written modulo $2^{\mathit{bits}}$ (its lowest
bits), except in format 8, which limits the change from one sample to the next. The formats
are listed by the library constant \code{WFDB\_FMT\_LIST}, in the order of
\cref{tab:formats}; every implementation reads and writes all of them.

\begin{xltabular}{\linewidth}{@{}rrP{5.4cm}Lc@{}}
\caption{The signal storage formats (bits per sample, the bytes of a group of samples, the
range of the values, and whether \wfdbx{} reads (R) and writes (W) them).}
\label{tab:formats}\\
\toprule
{\leavevmode\tabhead Format} & {\leavevmode\tabhead Bits} & {\leavevmode\tabhead Layout} & {\leavevmode\tabhead Range (invalid value)} &
  {\leavevmode\tabhead R/W} \\
\midrule
\endfirsthead
\toprule
{\leavevmode\tabhead Format} & {\leavevmode\tabhead Bits} & {\leavevmode\tabhead Layout} & {\leavevmode\tabhead Range (invalid value)} &
  {\leavevmode\tabhead R/W} \\
\midrule
\endhead
\bottomrule
\endfoot
0 & -- & no signal file (a null signal) & every sample invalid & R/W \\
8 & 8 & 1 byte: the difference from the previous sample, two's complement & differences
  $-128$ to 127 (no invalid value) & R/W \\
16 & 16 & 2 bytes, two's complement, little-endian & $-32768$ to 32767 ($-32768$) & R/W \\
61 & 16 & 2 bytes, two's complement, big-endian & $-32768$ to 32767 ($-32768$) & R/W \\
80 & 8 & 1 byte, offset binary (value + 128) & $-128$ to 127 ($-128$) & R/W \\
160 & 16 & 2 bytes, offset binary (value + 32768), little-endian & $-32768$ to 32767
  ($-32768$) & R/W \\
212 & 12 & 2 samples in 3 bytes & $-2048$ to 2047 ($-2048$) & R/W \\
310 & 10 & 3 samples in two 16-bit little-endian words (15 bits of each) & $-512$ to 511
  ($-512$) & R/W \\
311 & 10 & 3 samples in one 32-bit little-endian word (30 bits) & $-512$ to 511 ($-512$) &
  R/W \\
24 & 24 & 3 bytes, two's complement, little-endian & $-2^{23}$ to $2^{23}-1$ ($-2^{23}$) &
  R/W \\
32 & 32 & 4 bytes, two's complement, little-endian & $-2^{31}$ to $2^{31}-1$ ($-2^{31}$) &
  R/W \\
508 & 8 & FLAC stream, 8 bits per sample & $-128$ to 127 ($-128$) & R/W \\
516 & 16 & FLAC stream, 16 bits per sample & $-32768$ to 32767 ($-32768$) & R/W \\
524 & 24 & FLAC stream, 24 bits per sample & $-2^{23}$ to $2^{23}-1$ ($-2^{23}$) & R/W \\
\end{xltabular}

The ADC resolution given in the header (\cref{chap:records}) may be smaller than the
storage width of the format (11-bit samples in format 212, as in \code{100s}); the format
only has to hold the range of the values. The precise layout of each format follows. In
the bit diagrams, bit 0 is the least significant bit, and $s_0$, $s_1$, $s_2$ are
consecutive samples of the group.

\subsection{Format 0: null signals}
\label{sec:formats-0}

A signal of format 0 has no samples on disk (its file name in the header is usually
\code{\textasciitilde}). Every sample read is invalid (\code{WFDB\_INVALID\_SAMPLE}, or, in
\code{WFDB\_GVPAD} mode, the signal's initial value from the header); samples written to it
are discarded. Format 0 is used by the layout segment and the null segments of
multi-segment records.

\subsection{Format 8: first differences}
\label{sec:formats-8}

Each byte is the difference between a sample and the previous sample of the same signal,
as a signed 8-bit integer; the first difference is taken from the initial value given in
the header. The samples themselves can take any value, but they change by $-128$ to 127 from one
sample to the next: \code{putvec} limits a larger difference to $-128$ or 127 and returns 0 (the slew
rate is too high), and the following samples are then relative to the value
actually written. The format has no invalid value, and an error in the file affects every
following sample.

\subsection{Formats 16 and 61: 16-bit amplitudes}
\label{sec:formats-16}

Each sample is a 16-bit two's complement integer, least significant byte first (format 16)
or most significant byte first (format 61).

\subsection{Formats 80 and 160: offset binary}
\label{sec:formats-80}

Each sample is stored as an unsigned integer, the value plus $2^{\mathit{bits}-1}$: one byte,
the value plus 128 (format 80), or two bytes, the value plus 32768, least significant byte
first (format 160). Format 80 is that of 8-bit WAV files (\cmd{wav2mit}).

\subsection{Format 212: 12-bit amplitudes}
\label{sec:formats-212}

Two samples are packed into three bytes $b_0 b_1 b_2$:
\begin{itemize}
\item $s_0$: the low 8 bits in $b_0$, the high 4 bits in bits 0--3 of $b_1$;
\item $s_1$: the high 4 bits in bits 4--7 of $b_1$, the low 8 bits in $b_2$.
\end{itemize}
Each sample is a 12-bit two's complement integer. When the number of samples written is
odd, the last sample is written as the first two bytes of a group ($b_0 b_1$).

\subsection{Format 310: 10-bit amplitudes in 16-bit words}
\label{sec:formats-310}

Three samples are packed into two 16-bit little-endian words $w_0$ and $w_1$, bit 0 of each
word being unused (0):
\begin{itemize}
\item $s_0$: bits 1--10 of $w_0$;
\item $s_1$: bits 1--10 of $w_1$;
\item $s_2$: its low 5 bits in bits 11--15 of $w_0$, its high 5 bits in bits 11--15 of
  $w_1$.
\end{itemize}
Each sample is a 10-bit two's complement integer. An incomplete group at the end of the
file is written as the words that hold its samples ($w_0$ for one sample, $w_0 w_1$ for
two).

\subsection{Format 311: 10-bit amplitudes in 32-bit words}
\label{sec:formats-311}

Three samples are packed into one 32-bit little-endian word $w$ (four bytes), bits 30 and
31 being unused (0):
\begin{itemize}
\item $s_0$: bits 0--9 of $w$;
\item $s_1$: bits 10--19;
\item $s_2$: bits 20--29.
\end{itemize}
Each sample is a 10-bit two's complement integer. An incomplete group at the end of the
file is written as its first two bytes (one sample) or its four bytes (two samples).

\subsection{Formats 24 and 32: 24- and 32-bit amplitudes}
\label{sec:formats-24}

Each sample is a two's complement integer of 3 bytes (format 24) or 4 bytes (format 32),
least significant byte first. In format 32, the invalid value $-2^{31}$ differs from
\code{WFDB\_INVALID\_SAMPLE}: an invalid sample is written as $-2^{31}$, which is read back
as \code{WFDB\_INVALID\_SAMPLE}.

\section{FLAC formats}
\label{sec:formats-flac}

Formats 508, 516 and 524 store the samples of a signal group losslessly compressed, as a
FLAC stream (Free Lossless Audio Codec) of $\mathit{format} - 500$ bits per sample (8, 16
or 24). Each signal of the group is a channel of the stream, so a FLAC file holds at most 8
signals, and all of them must have the same number of samples per frame (the library
reports \code{cannot store N signals in a single FLAC file} or
\code{every signal in a FLAC file must have the same sampling frequency}). The samples of a
frame of a signal with several samples per frame are consecutive samples of its channel.
The FLAC sample rate is nominal (96\,kHz); the sampling frequency is the one given in the
header.

The C library reads and writes these formats with libFLAC. \wfdbx{} needs no external
library:
\begin{itemize}
\item \textbf{Reading} uses a translation of the stream decoder of libFLAC 1.4.3, in each
  implementation: every subframe type (constant, verbatim, fixed and linear prediction),
  stereo decorrelation, wasted bits, variable block sizes and CRC checks. Damaged files are
  decoded, resynchronized and reported as by the C library, and seeking
  (\code{isigsettime}) works as in C. The MD5 signature of the stream is checked when the
  input is closed if the environment variable \code{WFDB\_FLAC\_CHECK\_MD5} is set
  (\code{isigclose: warning: incorrect MD5 hash in FLAC input}).
\item \textbf{Writing} uses an encoder of \wfdbx{}'s own: the stream begins with
  \code{fLaC}, a \code{STREAMINFO} block (block size 4096, the frame sizes, the number of
  channels and bits per sample, the total number of samples and the MD5 signature, all
  completed when the file is closed) and an empty \code{VORBIS\_COMMENT} block (vendor
  \code{wfdb-x}). The samples are coded in frames of 4096 samples per channel (the last
  frame shorter), the channels independently, each as a constant, verbatim or
  fixed-predictor (order 0 to 4) subframe, whichever is smallest, with Rice-coded residuals
  (partition orders 0 to 6). A sample is first reduced to the width of the format.
\end{itemize}
The files written are standard FLAC streams: libFLAC, and therefore the C library, reads
them back with identical sample values. Their bytes differ from those that libFLAC writes
(which depend on its version and on the processor), so a record written in a FLAC format
by \wfdbx{} has the same samples and checksums as one written by the C library, but not the
same signal file.

\section{Annotation files}
\label{sec:formats-annotations}

An annotation file belongs to a record and an annotator: \code{100s.atr} holds the
annotations of the annotator \code{atr} for the record \code{100s}. Each annotation has a
time (a sample number, or a number of ticks of the file's time resolution), an annotation
code (\code{anntyp}, \cref{sec:formats-codes}), three small integer fields (\code{subtyp},
\code{chan} and \code{num}) and optional auxiliary information (\code{aux}, a string of up
to 255 bytes). The library reads two formats, MIT (the WFDB format, the default for
writing) and AHA; it recognizes the format of a file when it opens it for reading.

\subsection{MIT format}
\label{sec:formats-mit-annotations}

An MIT-format annotation file is a sequence of 16-bit words, each stored least significant
byte first. The 6 high bits of a word are a code $A$ (\code{(w >> 10) \& 63}), and the 10
low bits a value $I$ (\code{w \& 1023}). An annotation is an \emph{annotation word}, whose
code $A$ (1 to 49) is the annotation code and whose value $I$ is the interval, in ticks,
from the time of the previous annotation (from 0 for the first), followed by any number of
\emph{pseudo-annotation words} that modify it, with the codes 59 to 63 of
\cref{tab:mit-pseudo}. The 32-bit integers of the file are stored as two 16-bit words, the
most significant word first, each least significant byte first (the byte order of the
PDP-11: the bytes 2, 3, 0, 1 of a little-endian integer). The word 0 ($A = 0$, $I = 0$)
ends the file.

\begin{xltabular}{\linewidth}{@{}rlL@{}}
\caption{The pseudo-annotation codes of MIT-format annotation files.}
\label{tab:mit-pseudo}\\
\toprule
{\leavevmode\tabhead $A$} & {\leavevmode\tabhead Name} & {\leavevmode\tabhead Meaning} \\
\midrule
\endfirsthead
\toprule
{\leavevmode\tabhead $A$} & {\leavevmode\tabhead Name} & {\leavevmode\tabhead Meaning} \\
\midrule
\endhead
\bottomrule
\endfoot
59 & \code{SKIP} & $I$ is 0; the next two words are a signed 32-bit interval, added to
  the time. It precedes an annotation word whose interval does not fit in 10 bits (more
  than 1023 ticks, or a negative interval, in a file that is out of order), which then has
  the interval 0. Intervals beyond $\pm(2^{31}-1)$ take several \code{SKIP}s. \\
60 & \code{NUM} & $I$ is the \code{num} field of this annotation and of the following
  ones (written only when it changes; initially 0). \\
61 & \code{SUB} & $I$ is the \code{subtyp} field of this annotation only (written only when
  it is not 0). \\
62 & \code{CHN} & $I$ is the \code{chan} field of this annotation and of the following
  ones (written only when it changes; initially 0). \\
63 & \code{AUX} & The low 8 bits of $I$ are a byte count $n$; the next $n$ bytes are the
  \code{aux} string of this annotation, followed by a null byte if $n$ is odd, so that the
  file stays word-aligned. \\
\end{xltabular}

The \code{subtyp} and \code{num} fields are signed 8-bit integers and \code{chan} an
unsigned 8-bit integer, stored in the 10 bits of $I$. Codes 50 to 58 are not used. The
codes 59 to 63 are reserved for the pseudo-annotation words, and cannot be annotation
codes: an annotation word with one of them would be read as a control word, and would
corrupt the rest of the file. The times \code{WFDB\_TIME\_MIN} and \code{WFDB\_TIME\_MAX}
($-2^{63}$ and $2^{63}-1$) cannot be written either. The library function \code{putann}
refuses these times (\code{time overflow}) but does not check the code, as in C (a code
above 63 also loses its high bits); the high-level \code{write\_annotations} accepts only
the codes 0 to 58 and refuses the others, and these times, before it writes anything
(\cref{sec:additions}). An annotation of code 0 (\code{NOTQRS}), which a word of value 0 could not represent, is
written as a \code{SKIP} of its interval minus 1 followed by a word of code 0 and interval
1.

At the beginning of a file, \code{NOTE} annotations at time 0 (with \code{subtyp},
\code{chan} and \code{num} 0) whose \code{aux} begins with \code{\#\#} carry information
for the library, written by \code{putann} before the first annotation and read by
\code{annopen}:
\begin{itemize}
\item \code{\#\# time resolution: }\textit{F}: the times of the file are in ticks of
  \textit{F} per second, not in sample intervals of the record (written when
  \code{setafreq} has set a resolution, as \cmd{wrann} does with the sampling frequency of
  the record; \code{getiaorigfreq} returns it);
\item \code{\#\# annotation type definitions}, then one annotation per code whose mnemonic
  or description has been redefined (by \code{setannstr} or \code{setanndesc} with a
  positive code), with the \code{aux} \textit{code mnemonic description}, then
  \code{\#\# end of definitions}. When it reads such a file, the library takes these
  definitions as the mnemonics and descriptions of \code{annstr}, \code{strann} and
  \code{anndesc};
\item after these, an annotation of code 0 at time 0 ends the table.
\end{itemize}

An MIT-format file written by \code{putann} need not be in canonical order (by time, then
\code{num}, then \code{chan}); \wfdbx{} sorts an out-of-order file itself when it is closed,
without losing annotations (\cref{chap:behavior}), unless the environment variable
\code{WFDBANNSORT} is 0.

\subsection{AHA format}
\label{sec:formats-aha-annotations}

An AHA-format annotation file, as used by the AHA Database and written for it by
\cmd{m2a}, has 16 bytes per annotation:

\begin{xltabular}{\linewidth}{@{}lL@{}}
\caption{An annotation of an AHA-format annotation file.}
\label{tab:aha-annotation}\\
\toprule
{\leavevmode\tabhead Bytes} & {\leavevmode\tabhead Contents} \\
\midrule
\endfirsthead
\toprule
{\leavevmode\tabhead Bytes} & {\leavevmode\tabhead Contents} \\
\midrule
\endhead
\bottomrule
\endfoot
0 & 0 \\
1 & the AHA annotation code, a character (\cref{tab:aha-codes}); mapped to an annotation
  code by \code{ammap} when read, and from one by \code{mamap} when written \\
2--5 & the time, a 32-bit integer in PDP-11 order (most significant word first) \\
6--7 & the serial number of the annotation, a 16-bit little-endian integer, from 1 \\
8 & the \code{subtyp} field \\
9 & the annotation code (\code{anntyp}) when written; read as the \code{chan} field \\
10--15 & the \code{aux} string, up to 6 bytes, padded with null bytes \\
\end{xltabular}

The file is padded with bytes 255 (\code{0377}) to a multiple of 1024 bytes; a byte 255 in
place of the first byte of an annotation ends it. When it reads an AHA file, the library
gives an AHA code \code{U} with a \code{subtyp} of 0 the \code{subtyp} $-1$, and
\code{mamap} maps \code{NOISE} to \code{U} only with that subtype. A file opened for
reading is taken to be in AHA format when its first byte is 0 and its second byte is an AHA
code other than \code{[} and \code{]}; the library warns (\code{file appears to be in AHA
format}) when the format it finds is not the one requested. The times are in the units
written: sample intervals for files written by \cmd{m2a} without \opt{-s}, and milliseconds
in the AHA Database tape files (\cmd{a2m} \opt{-t 1} and \opt{-t 2}). The compressed
annotation files of the AHA Database CD-ROMs (\file{.ANO}), read only by \cmd{a2m}
(\opt{-t 3}), have 6 bytes per annotation: the AHA code, a 24-bit time in milliseconds (most
significant byte first; 255 in its high byte marks the end), and a 16-bit serial number.

\section{Annotation codes}
\label{sec:formats-codes}

Annotation codes are integers from 1 to \code{ACMAX} (49); 0, \code{NOTQRS}, is not an
annotation code (\code{isann} is false for it), and codes 42 to 49 have no standard meaning
(they can be defined by \code{setannstr} and \code{setanndesc}, and by the definitions of
an annotation file). Each code has a symbolic name (a constant of the library: \code{NORMAL},
\code{PVC}, ...), a mnemonic (the text of \cmd{rdann}, returned by \code{annstr} and
\code{ecgstr}) and a description (\code{anndesc}). \Cref{tab:anncodes} gives the standard
values, those of the library when no annotation file has redefined them.

\begin{xltabular}{\linewidth}{@{}rP{3.3cm}lL@{}}
\caption{The annotation codes: number, symbolic name, mnemonic and description.}
\label{tab:anncodes}\\
\toprule
{\leavevmode\tabhead Code} & {\leavevmode\tabhead Name} & {\leavevmode\tabhead Mnemonic} & {\leavevmode\tabhead Description} \\
\midrule
\endfirsthead
\toprule
{\leavevmode\tabhead Code} & {\leavevmode\tabhead Name} & {\leavevmode\tabhead Mnemonic} & {\leavevmode\tabhead Description} \\
\midrule
\endhead
\bottomrule
\endfoot
0 & \code{NOTQRS} & (space) & (not an annotation code) \\
1 & \code{NORMAL} & \code{N} & Normal beat \\
2 & \code{LBBB} & \code{L} & Left bundle branch block beat \\
3 & \code{RBBB} & \code{R} & Right bundle branch block beat \\
4 & \code{ABERR} & \code{a} & Aberrated atrial premature beat \\
5 & \code{PVC} & \code{V} & Premature ventricular contraction \\
6 & \code{FUSION} & \code{F} & Fusion of ventricular and normal beat \\
7 & \code{NPC} & \code{J} & Nodal (junctional) premature beat \\
8 & \code{APC} & \code{A} & Atrial premature beat \\
9 & \code{SVPB} & \code{S} & Supraventricular premature or ectopic beat \\
10 & \code{VESC} & \code{E} & Ventricular escape beat \\
11 & \code{NESC} & \code{j} & Nodal (junctional) escape beat \\
12 & \code{PACE} & \code{/} & Paced beat \\
13 & \code{UNKNOWN} & \code{Q} & Unclassifiable beat \\
14 & \code{NOISE} & \code{\textasciitilde{}} & Change in signal quality \\
15 & -- & \code{[15]} & (none) \\
16 & \code{ARFCT} & \code{\textbar{}} & Isolated QRS-like artifact \\
17 & -- & \code{[17]} & (none) \\
18 & \code{STCH} & \code{s} & ST segment change \\
19 & \code{TCH} & \code{T} & T-wave change \\
20 & \code{SYSTOLE} & \code{*} & Systole \\
21 & \code{DIASTOLE} & \code{D} & Diastole \\
22 & \code{NOTE} & \code{\textquotedbl{}} & Comment annotation \\
23 & \code{MEASURE} & \code{=} & Measurement annotation \\
24 & \code{PWAVE} & \code{p} & P-wave peak \\
25 & \code{BBB} & \code{B} & Bundle branch block beat (unspecified) \\
26 & \code{PACESP} & \code{\textasciicircum{}} & (Non-captured) pacemaker artifact \\
27 & \code{TWAVE} & \code{t} & T-wave peak \\
28 & \code{RHYTHM} & \code{+} & Rhythm change \\
29 & \code{UWAVE} & \code{u} & U-wave peak \\
30 & \code{LEARN} & \code{?} & Beat not classified during learning \\
31 & \code{FLWAV} & \code{!} & Ventricular flutter wave \\
32 & \code{VFON} & \code{[} & Start of ventricular flutter/fibrillation \\
33 & \code{VFOFF} & \code{]} & End of ventricular flutter/fibrillation \\
34 & \code{AESC} & \code{e} & Atrial escape beat \\
35 & \code{SVESC} & \code{n} & Supraventricular escape beat \\
36 & \code{LINK} & \code{@} & Link to external data (aux contains URL) \\
37 & \code{NAPC} & \code{x} & Non-conducted P-wave (blocked APC) \\
38 & \code{PFUS} & \code{f} & Fusion of paced and normal beat \\
39 & \code{WFON}, \code{PQ} & \code{(} & Waveform onset \\
40 & \code{WFOFF}, \code{JPT} & \code{)} & Waveform end \\
41 & \code{RONT} & \code{r} & R-on-T premature ventricular contraction \\
42 & -- & \code{[42]} & (none) \\
43 & -- & \code{[43]} & (none) \\
44 & -- & \code{[44]} & (none) \\
45 & -- & \code{[45]} & (none) \\
46 & -- & \code{[46]} & (none) \\
47 & -- & \code{[47]} & (none) \\
48 & -- & \code{[48]} & (none) \\
49 & -- & \code{[49]} & (none) \\
\end{xltabular}

\subsection{Mapping tables}
\label{sec:formats-maps}

The library classifies the annotation codes with tables that programs can query and change
(\cref{chap:library}): \code{isqrs(a)} is 1 for the codes of QRS complexes (beats);
\code{map1(a)} maps a code to one of \code{NOTQRS}, \code{NORMAL}, \code{PVC},
\code{FUSION} and \code{LEARN}; \code{map2(a)} to one of those or \code{SVPB}; \code{mamap(a, s)} gives the AHA
code (a character) of an annotation code with the subtype \code{s}; and \code{annpos(a)}
the position at which an annotation is drawn (\code{APSTD} at the signal, \code{APHIGH}
above it, \code{APLOW} below, \code{APUNDEF} undefined; the position codes \code{APATT} (4,
attached to the signal given by \code{chan}), \code{APAHIGH} (5) and \code{APALOW} (6, above
and below it) are defined but used by no standard code).
\code{setisqrs}, \code{setmap1}, \code{setmap2} and \code{setannpos} change the tables. A
code out of range gives 0, \code{NOTQRS}, \code{NOTQRS}, \code{O} and \code{APUNDEF}.
\Cref{tab:annmaps} gives the tables for the codes 0 to 41 (codes 42 to 49 give 0,
\code{NOTQRS}, \code{NOTQRS}, \code{O} and \code{APUNDEF}).

\begin{xltabular}{\linewidth}{@{}rlrlllL@{}}
\caption{The mapping tables of the annotation codes 0 to 41 (\code{mamap} of code 14 is
\code{U} with the subtype $-1$, \code{O} otherwise).}
\label{tab:annmaps}\\
\toprule
{\leavevmode\tabhead Code} & {\leavevmode\tabhead Mnemonic} & {\leavevmode\tabhead isqrs} & {\leavevmode\tabhead map1} & {\leavevmode\tabhead map2} & {\leavevmode\tabhead mamap} & {\leavevmode\tabhead annpos} \\
\midrule
\endfirsthead
\toprule
{\leavevmode\tabhead Code} & {\leavevmode\tabhead Mnemonic} & {\leavevmode\tabhead isqrs} & {\leavevmode\tabhead map1} & {\leavevmode\tabhead map2} & {\leavevmode\tabhead mamap} & {\leavevmode\tabhead annpos} \\
\midrule
\endhead
\bottomrule
\endfoot
0 & \code{\ } & 0 & \code{NOTQRS} & \code{NOTQRS} & \code{O} & \code{APUNDEF} \\
1 & \code{N} & 1 & \code{NORMAL} & \code{NORMAL} & \code{N} & \code{APSTD} \\
2 & \code{L} & 1 & \code{NORMAL} & \code{NORMAL} & \code{N} & \code{APSTD} \\
3 & \code{R} & 1 & \code{NORMAL} & \code{NORMAL} & \code{N} & \code{APSTD} \\
4 & \code{a} & 1 & \code{NORMAL} & \code{SVPB} & \code{N} & \code{APSTD} \\
5 & \code{V} & 1 & \code{PVC} & \code{PVC} & \code{V} & \code{APSTD} \\
6 & \code{F} & 1 & \code{FUSION} & \code{FUSION} & \code{F} & \code{APSTD} \\
7 & \code{J} & 1 & \code{NORMAL} & \code{SVPB} & \code{N} & \code{APSTD} \\
8 & \code{A} & 1 & \code{NORMAL} & \code{SVPB} & \code{N} & \code{APSTD} \\
9 & \code{S} & 1 & \code{NORMAL} & \code{SVPB} & \code{N} & \code{APSTD} \\
10 & \code{E} & 1 & \code{PVC} & \code{PVC} & \code{E} & \code{APSTD} \\
11 & \code{j} & 1 & \code{NORMAL} & \code{NORMAL} & \code{N} & \code{APSTD} \\
12 & \code{/} & 1 & \code{NORMAL} & \code{NORMAL} & \code{P} & \code{APSTD} \\
13 & \code{Q} & 1 & \code{NORMAL} & \code{NORMAL} & \code{Q} & \code{APSTD} \\
14 & \code{\textasciitilde{}} & 0 & \code{NOTQRS} & \code{NOTQRS} & \code{U} & \code{APHIGH} \\
15 & \code{[15]} & 0 & \code{NOTQRS} & \code{NOTQRS} & \code{O} & \code{APUNDEF} \\
16 & \code{\textbar{}} & 0 & \code{NOTQRS} & \code{NOTQRS} & \code{O} & \code{APHIGH} \\
17 & \code{[17]} & 0 & \code{NOTQRS} & \code{NOTQRS} & \code{O} & \code{APUNDEF} \\
18 & \code{s} & 0 & \code{NOTQRS} & \code{NOTQRS} & \code{O} & \code{APHIGH} \\
19 & \code{T} & 0 & \code{NOTQRS} & \code{NOTQRS} & \code{O} & \code{APHIGH} \\
20 & \code{*} & 0 & \code{NOTQRS} & \code{NOTQRS} & \code{O} & \code{APHIGH} \\
21 & \code{D} & 0 & \code{NOTQRS} & \code{NOTQRS} & \code{O} & \code{APHIGH} \\
22 & \code{\textquotedbl{}} & 0 & \code{NOTQRS} & \code{NOTQRS} & \code{O} & \code{APHIGH} \\
23 & \code{=} & 0 & \code{NOTQRS} & \code{NOTQRS} & \code{O} & \code{APHIGH} \\
24 & \code{p} & 0 & \code{NOTQRS} & \code{NOTQRS} & \code{O} & \code{APHIGH} \\
25 & \code{B} & 1 & \code{NORMAL} & \code{NORMAL} & \code{N} & \code{APSTD} \\
26 & \code{\textasciicircum{}} & 0 & \code{NOTQRS} & \code{NOTQRS} & \code{O} & \code{APHIGH} \\
27 & \code{t} & 0 & \code{NOTQRS} & \code{NOTQRS} & \code{O} & \code{APHIGH} \\
28 & \code{+} & 0 & \code{NOTQRS} & \code{NOTQRS} & \code{O} & \code{APLOW} \\
29 & \code{u} & 0 & \code{NOTQRS} & \code{NOTQRS} & \code{O} & \code{APHIGH} \\
30 & \code{?} & 1 & \code{LEARN} & \code{LEARN} & \code{Q} & \code{APSTD} \\
31 & \code{!} & 1 & \code{PVC} & \code{PVC} & \code{O} & \code{APSTD} \\
32 & \code{[} & 0 & \code{NOTQRS} & \code{NOTQRS} & \code{[} & \code{APSTD} \\
33 & \code{]} & 0 & \code{NOTQRS} & \code{NOTQRS} & \code{]} & \code{APSTD} \\
34 & \code{e} & 1 & \code{NORMAL} & \code{NORMAL} & \code{N} & \code{APSTD} \\
35 & \code{n} & 1 & \code{NORMAL} & \code{NORMAL} & \code{N} & \code{APSTD} \\
36 & \code{@} & 0 & \code{NOTQRS} & \code{NOTQRS} & \code{O} & \code{APHIGH} \\
37 & \code{x} & 0 & \code{NOTQRS} & \code{NOTQRS} & \code{O} & \code{APHIGH} \\
38 & \code{f} & 1 & \code{NORMAL} & \code{NORMAL} & \code{N} & \code{APSTD} \\
39 & \code{(} & 0 & \code{NOTQRS} & \code{NOTQRS} & \code{O} & \code{APHIGH} \\
40 & \code{)} & 0 & \code{NOTQRS} & \code{NOTQRS} & \code{O} & \code{APHIGH} \\
41 & \code{r} & 1 & \code{PVC} & \code{PVC} & \code{R} & \code{APSTD} \\
\end{xltabular}

\code{ammap(c)}, the reverse of \code{mamap}, maps an AHA code (the character \code{c},
from \code{E} to \code{]}) to an annotation code, and gives \code{NOTQRS} for any other
character:

\begin{xltabular}{\linewidth}{@{}lL@{}}
\caption{The AHA annotation codes and the annotation codes that \code{ammap} gives them.}
\label{tab:aha-codes}\\
\toprule
{\leavevmode\tabhead AHA code} & {\leavevmode\tabhead ammap} \\
\midrule
\endfirsthead
\toprule
{\leavevmode\tabhead AHA code} & {\leavevmode\tabhead ammap} \\
\midrule
\endhead
\bottomrule
\endfoot
\code{E} & \code{VESC} \\
\code{F} & \code{FUSION} \\
\code{G} & \code{NOTQRS} \\
\code{H} & \code{NOTQRS} \\
\code{I} & \code{NOTQRS} \\
\code{J} & \code{NOTQRS} \\
\code{K} & \code{NOTQRS} \\
\code{L} & \code{NOTQRS} \\
\code{M} & \code{NOTQRS} \\
\code{N} & \code{NORMAL} \\
\code{O} & \code{NOTE} \\
\code{P} & \code{PACE} \\
\code{Q} & \code{UNKNOWN} \\
\code{R} & \code{RONT} \\
\code{S} & \code{NOTQRS} \\
\code{T} & \code{NOTQRS} \\
\code{U} & \code{NOISE} \\
\code{V} & \code{PVC} \\
\code{W} & \code{NOTQRS} \\
\code{X} & \code{NOTQRS} \\
\code{Y} & \code{NOTQRS} \\
\code{Z} & \code{NOTQRS} \\
\code{[} & \code{VFON} \\
\code{\textbackslash{}} & \code{NOTQRS} \\
\code{]} & \code{VFOFF} \\
\end{xltabular}
